\documentclass[11pt,a4paper]{article}
\usepackage[margin=25mm]{geometry}
\usepackage[T1]{fontenc}
\usepackage[utf8]{inputenc}
\usepackage{lmodern,microtype,amsmath,amssymb,booktabs,longtable,array}
\usepackage{graphicx,xcolor,enumitem}
\usepackage[numbers,sort&compress]{natbib}
\usepackage[colorlinks=true,linkcolor=blue!45!black,citecolor=blue!45!black,urlcolor=blue!45!black]{hyperref}
\usepackage{listings}
\lstset{basicstyle=\ttfamily\small,breaklines=true,columns=fullflexible,frame=single,framerule=0.2pt}
\setlist{nosep,leftmargin=*}
\setlength{\emergencystretch}{2em}
\newcommand{\vt}{v_{\rm t}}
\newcommand{\smu}{S\mu^2}
\newcommand{\cm}{\,\mathrm{cm}^{-1}}
\title{Methanol L33T\\\large Molecular data, construction and reproducibility}
\date{20 November 2026\\\small Release v1.0 \quad Molecular model \texttt{l33t-v3}}
\begin{document}
\maketitle
\begin{abstract}
L33T provides separate A- and E-symmetry molecular files for ordinary and
carbon-13 methanol, with a large torsional and vibrational excitation basis.
The source package reconstructs the torsion--rotation spectrum, electric-dipole
strengths, Einstein coefficients and a source-based collision network, including
the specified extrapolations. This document defines the numerical model and its
reproducibility contract. L33T v3 is the adopted molecular dataset for numerical
experiments, selected by the maintainer on 6 November 2026. Its collision and
vibrational prescriptions include the approximations specified here; their
physical accuracy must be assessed for each application. Public release v1.0 supplies these four molecular files, their row-level
provenance and the reconstruction code.
\end{abstract}
\tableofcontents
\clearpage

\section{Products and intended use}
The package targets excitation calculations that require high-lying methanol
levels and infrared pumping channels. A and E species are separate systems;
the files contain no A--E conversion. The four products are listed in
Table~\ref{tab:products}. Their identities are fixed by SHA-256 checksums in
\path{provenance/products.json}, including the three accompanying provenance
tables for each file. Filenames alone are insufficient version identifiers.

\begin{table}[ht]
\centering
\caption{Frozen L33T v3 products. Filenames end in
\texttt{-l33t-E2000-J46.dat}. All collision rows are downward coefficients.}
\label{tab:products}
\begin{tabular}{lrr}
\toprule File prefix & Levels & Radiative rows\\\midrule
\texttt{a-12ch3oh} & 7,600 & 146,635\\
\texttt{e-12ch3oh} & 7,669 & 163,932\\
\texttt{a-13ch3oh} & 7,714 & 149,820\\
\texttt{e-13ch3oh} & 7,794 & 167,491\\\midrule
Total & 30,777 & 627,878\\\bottomrule
\end{tabular}
\end{table}

Together the products contain 7,811,683 collision rows. Each retains model
states with $E<2000\cm$ above the ground state of its own symmetry species.
The native enumeration uses $J\le46$ and $|K|\le22$; the energy cut removes
many states in this box. The retained manifolds are $\vt=0,\ldots,4$ and six
uncoupled vibrational fundamentals: CO stretch ($\nu_8$), in-plane CH$_3$
rock ($\nu_7$), $\nu_6$, $\nu_5$, $\nu_4$ and $\nu_{10}$.
This is completeness within the declared model and enumeration, not all
physical rovibrational states below the energy threshold. For example, the
weak $\nu_{11}$ band is omitted, and coupled combination states are not
constructed.

The intended exploratory ranges were $T_{\rm kin}\le400$ K and
$T_{\rm dust}\le500$ K. They are model-design ranges, not a demonstrated
accuracy domain. The tabulated collision grid is
\begin{equation}
T/\mathrm K=(10,20,30,50,75,100,150,200,300,400).
\end{equation}
Published rate calculations do not cover this entire level--temperature space.
The approximation labels in the collision ledger must accompany any use of
the high-temperature or high-$J$ model.

\section{Hamiltonian, symmetry and state identity}
\subsection{Frozen torsion--rotation model}
The main-isotope calculation uses the fitted methanol Hamiltonian of
Xu et al.~\cite{XuEtAl2008}; the isotope calculation uses the separate
$^{13}$CH$_3$OH parameter set of Xu et al.~\cite{XuEtAl2014C13}.
The implementation is a Python interface to the checksum-identified native
BELGI Fortran matrix routines and supplied prediction drivers. Python stages
coefficients and compiles the native calculation. BELGI supplies Hamiltonian
matrices; local overlays perform the A-species parity-block eigensolve, state
assignment and all cached line-strength contractions. The GLOBALFIT program
is not compiled or executed. Native unpruned INTEN is used only in the separate
contraction comparison, which shares the repaired eigensystem. Users should acknowledge BELGI-C$_s$ and cite
Hougen, Kleiner and Godefroid~\cite{Hougen_1994}, in addition to the
isotope-specific parameter sources. No refit is performed.

A schematic leading-order form, using the positive coupling sign in the
archived native torsional basis, is
\begin{align}
H={}&F(p_\gamma+\rho J_a)^2+
\sum_{n\ge1}\frac{V_{3n}}{2}[1-\cos(3n\gamma)]\nonumber\\
&+A_{\rm RAM}J_a^2+B_{\rm RAM}J_b^2+C_{\rm RAM}J_c^2
+D_{ab}\{J_a,J_b\}+H_{\rm higher}.
\label{eq:ham}
\end{align}
Here $\gamma$ is the internal-rotation angle, angular-momentum operators are
dimensionless, and the coefficients carry the corresponding spectroscopic
units. Equation~\eqref{eq:ham} describes the operator structure; it is not a
replacement for the full fitted operator inventory. The native torsional
factor is $3m+\sigma+\rho K$. Changing the sign convention for $\gamma$ or
the body axes requires a consistent parameter and dipole transformation.

The main-isotope loader retains all 119 nonzero fitted coefficients and
explicitly fills the other 86 slots in the 205-entry constant vector with
zero. The isotope has its own frozen vector. Native matrices, fitted
coefficients and patch anchors are checksum-bound; no unspecified compiler
initialization supplies missing coefficients. The fitted $\vt\le2$ domain
does not make $\vt=3,4$ or the high-$J$ extension independently validated
spectroscopy.

\subsection{Parity and complete eigenvectors}
Nearly degenerate A-state pairs require a parity-resolved eigensolve.
The overlay first diagonalizes the physical parity operator, then solves
the Hamiltonian within its separate parity subspaces. It checks the parity
operator, Hamiltonian commutator, eigen-residuals and orthogonality. A global
bijective overlap assignment maintains the exported A identities across $J$.
The dipole contraction retains all nine internal blocks of the native
torsion--rotation vectors. The E calculation retains the native ENCAL2
assignment and the same full-vector contraction.

Labels have the form \texttt{A:vt0:J5:K2:+} or
\texttt{E:vt0:J5:K-2}. A uses $K\ge0$ and a separate sign; E uses signed
$K$. The A label sign is not itself physical parity: for torsional states,
\begin{equation}
P_A=s_A(-1)^{J+\vt}.
\end{equation}
The $\nu_{10}$ level retains its carrier ground-state label but has the
opposite physical parity. Assignment is a reproducible convention in a
mixed-state system; it does not independently prove every approximate
spectroscopic $\vt,K$ label at a resonance.

Energies are referred to the ground of each species and weights are
$g=2J+1$. No common nuclear-spin multiplier is inserted in these weights.
For deterministic serialization only, exact level-energy ties are separated
with \texttt{nextafter}; the level ledger retains both physical and serialized
energies. Artificial tie splittings do not generate physical radiative lines.

\section{Electric-dipole strengths and Einstein coefficients}
\subsection{Native operators}
The permanent torsion--rotation operator uses the adopted constant components
$\mu_a=0.899$ D and $\mu_b=-1.44$ D, with the three
$\cos3\gamma$ dipole coefficients fixed to zero. It is a specified
constant-dipole model, not a recovered general torsion-dependent dipole surface.
Rotational tensors are contracted with the complete mixed eigenvectors;
nominal $\Delta K$ filters are not imposed on the resulting amplitudes.

The vibrational construction reuses ground-manifold eigenfunctions with
effective band moments. A pure $a$-axis reference operator of 0.27 D supplies
CO, rock and the additional $A'$ bands. A direct $c$-axis reference supplies
$\nu_{10}$. In the adopted $a$-quantized basis,
\begin{equation}
\mu_b\ \propto\frac{T_{-1}-T_{+1}}{\sqrt2},\qquad
\mu_c\ \propto\frac{i(T_{-1}+T_{+1})}{\sqrt2}.
\end{equation}
The common phase $i$ cancels from the squared amplitude, but the relative
sign does not. The $c$ contraction therefore cannot generally be replaced
by a relabelled $b$ contraction. It includes transitions between the $\vt=0$
function carrying the vibrational level and lower torsional functions through
$\vt=4$.

\subsection{Units, frequency hierarchy and cuts}
For frequency in MHz and $\smu$ in D$^2$, the preserved conversion is
\begin{equation}
A_{ul}[\mathrm{s}^{-1}]
=1.16395\times10^{-20}\,
\frac{\nu_{ul}[\mathrm{MHz}]^3\,\smu[\mathrm D^2]}{2J_u+1}.
\label{eq:einstein}
\end{equation}
The numerical prefactor is the frozen implementation's rounded convention.
Strengths and weights must share that convention: catalogue nuclear-spin
factors cannot be imported without conversion. Frequencies and upper-state
energies are distinct output quantities.

The default frequency is the physical model gap, with
$\nu[\mathrm{MHz}]=29979.2458\,\Delta E[\mathrm{cm}^{-1}]$.
For main-isotope torsional lines, the frozen accepted Xu observation map
replaces the model frequency by exact endpoint identity. For the isotope,
negative-tag observed CDMS 033502 entries are selected by exact identity.
Explicit configured frequency overrides have final precedence. Ambiguous
identities are not guessed. The Einstein coefficient is recalculated with
the adopted frequency cubed, preserving the underlying strength.

The permanent native spectrum is generated down to $10^{-6}$ D$^2$.
The 0.27-D reference operators use that same native strength floor before
band-moment scaling. At product assembly, ordinary radiative rows require
$\smu\ge10^{-4}$ D$^2$, whereas $\nu_{10}$ cross-band rows require
$\smu\ge10^{-7}$ D$^2$. These are strength cuts, not fixed $A$ cuts.
The additional weak $\nu_{10}$ branches distinguish v3 from v2.
Zero-gap emission, duplicate endpoints, $|\Delta J|>1$, $J=0\to0$ and
forbidden A-parity transitions are rejected or excluded as specified by the
builder. Radiative graphs must remain connected.

\section{Vibrational energies and intensities}
CO and rock empirical terms derive from the archived Xu et al. V7/V8 line
compilation~\cite{Xu_2004}. Lower terms are bound to the ground identities
at $J\le22$, with a $2.5\times10^{-4}\cm$ tolerance. This avoids silently
extending an empirical identity match into the less accurate high-$J$
Hamiltonian prediction. Missing CO levels use the constant offset of the
highest measured $J$ in the same species/$K$/parity ladder, where available,
and otherwise a species-median offset. The rock policy interpolates offsets
in $J(J+1)$, keeps a constant offset above the last measured $J$, and uses
the declared $K$-origin proxies where needed.

The CO operator is normalized against thirteen specified published
absorption features at 297 K. The inverse median predicted/measured ratio
sets one common strength scale. One unresolved high-$K$ feature is excluded.
This is an in-sample band normalization, not an independent intensity test.
The isotope CO energies use the Mukhopadhyay et al. subband model
\cite{Mukhopadhyay_1995}; rock origins use the retained Predoi-Cross
transcription\cite{PredoiCross1997}. Main-isotope CO calibration and selected
rock strengths are transferred to the isotope by the specified identities.

For the four new bands, the carrier ground rotational energy is shifted by
the species origin in Table~\ref{tab:bands}. The listed E origins use the
common A zero of the source table; the builder subtracts its isotope-specific
A--E ground splitting before adding an E-species ground term. Origins and
moments follow the adopted values from Sunaga et al.~\cite{Sunaga_2026}.
The $^{13}$C files use the same four origins without an isotope shift: this
is a labelled proxy, not an isotope measurement.

\begin{table}[ht]
\centering\small
\caption{Additional band inputs. Origins are in cm$^{-1}$ and squared moments
in $e^2a_0^2$; multiply the latter by $2.541746473^2$ for D$^2$. Assigning
all $A'$ strength to $a$ is a model choice, not a unique consequence of symmetry.}
\label{tab:bands}
\begin{tabular}{lrrrrc}
\toprule Band & Origin A & Origin E & $\mu_A^2$ & $\mu_E^2$ & Axis\\\midrule
$\nu_6$ &1320.6&1297.5&$6.14\,10^{-4}$&$2.42\,10^{-4}$&$a$\\
$\nu_5$ &1453.3&1462.1&$1.48\,10^{-4}$&$1.41\,10^{-4}$&$a$\\
$\nu_4$ &1486.1&1483.3&$1.74\,10^{-4}$&$1.31\,10^{-4}$&$a$\\
$\nu_{10}$ &1481.5&1473.9&$1.29\,10^{-4}$&$9.67\,10^{-5}$&$c$\\\bottomrule
\end{tabular}
\end{table}

Intraband radiation copies the ground-state permanent strengths by exact
rotational identity and recomputes $A$ with the new energy gaps. Corrected
direct V7 strengths override the theoretical parallel-rock strengths for
their verified endpoints. The quarantined ambiguous source row remains
excluded. The model does not solve a coupled vibrational Hamiltonian or
include general rovibrational intensity borrowing.

\section{Collision data and extrapolation}
\subsection{Normalization and precedence}
The sources are the Rabli--Flower para-H$_2$ and He calculations for
$\vt=0,1,2$\cite{RabliFlower2010H2,RabliFlower2010He}, the torsion-changing
He calculation\cite{RabliFlower2011}, and the lower-level Dagdigian
para-/ortho-H$_2$ matrices\cite{Dagdigian2024}. The retained BASECOL releases
provide 18 matrices. Normalization joins each rate to source-level quantum
identities, energies and raw hashes. Of 425,430 raw rows, 6,108 have an endpoint
outside the historical collision identity basis; the remaining 419,322 rows
reproduce the frozen normalized master. Enlarging this basis is a scientific
change, not a packaging correction.

The E $\vt=1$ para-H$_2$ table IDC474 has a recorded discontinuity in its
180--200 K columns. These columns are retained in the raw file but excluded
from usable normalization. The 170 K endpoint therefore controls its
high-temperature continuation. Other tables retain their own usable grids.

Direct calculated rows take precedence where available; Dagdigian supplies
the low-level H$_2$ coverage. Native $\vt=1,2$ rows are retained. Missing
within-manifold rows in excited torsional and vibrational systems receive
exact-identity copies of ground-state coefficients. Native He torsion-changing
rows are retained, but new intermanifold couplings are not manufactured.
There is no broad high-$J$ ortho-H$_2$ extrapolation.

\subsection{The E-species hybrid}
The adopted E prescription is an explicit approximation:
\begin{equation}
k_{E,c}^{\rm hybrid}(T)=k_{E,p}^{\rm RF}(T)\,R_{A,c}(T),
\qquad c=p,o.
\end{equation}
The multiplier compares Dagdigian and Rabli--Flower A rates after summing
the parity partners of the corresponding $(J,|K|)$ pair. Its direction-neutral
summands are
\begin{equation}
\widetilde k_{ul}=g_u k_{ul}
 \exp[-c_2(E_u-E_l)/(2T)],\qquad
c_2=1.4387768775039338\ \mathrm{K\,cm}.
\end{equation}
The denominator uses RF para-H$_2$ for both multipliers. Missing quartet
ratios fall back to the median at the same $(|\Delta J|,|\Delta|K||)$,
then the global median. Rates are first oriented and interpolated as in the
implementation. The hybrid preserves RF E propensities while transferring
an A-based magnitude correction; it is not a new E scattering result.
An RF E zero is inherited even when the replaced Dagdigian coefficient is
positive: the locked inputs contain 370 such pair--temperature nodes per
collider, or 37 pairs at all ten temperatures. With identity copies these yield
333 all-zero para-H$_2$ rows and 407 all-zero ortho-H$_2$ rows in each E file.
Across all sources and extrapolations, E-12 and E-13 contain 55,935 and 57,133
all-zero para-H$_2$ rows; E-12 includes 49,462 extrapolated rows of this kind.
These are represented zeros, distinct from absent rows.
Exact zero quartet ratios are valid and do not trigger a fallback.
There are no zero quartet ratios in the locked inputs; the corrected
zero-handling rule therefore leaves the frozen products unchanged.

\subsection{Signed J/K donor rule}
Missing intramanifold para-H$_2$ pairs with $|\Delta J|\le10$ are eligible
for extrapolation. The donor edge is $J_d=12$. Existing rates are never
replaced. The first search preserves the two $K$ labels and A signs, shifting
both $J$ values down to the edge. If that exact jump is unavailable, clamped
or edge-anchored donors with a smaller, nonzero, equally signed $\Delta J$
may be used. Their coefficient is multiplied by
$|\Delta J_d|/|\Delta J_{\rm target}|$. A zero-$\Delta J$ target requires
a zero-$\Delta J$ donor and uses scale one.

If this search fails, define
\begin{equation}
s=\max(|K_u|,|K_l|)-(12-|J_u-J_l|).
\end{equation}
For $s>0$, both endpoints undergo
$J_i\mapsto J_i-s$ and $K_i\mapsto K_i-s\sigma_K$, where $\sigma_K$
is the sign of the endpoint with larger $|K|$ (upper endpoint on a tie).
The same donor search follows. Invalid $|K|>J$, negative A $K$, absent
donors and inappropriate branches remain unsupported. E $K$ is signed:
the translation can change the individual sign class while preserving signed
$\Delta K$. A $K=0$ uses its canonical $+$ label. Every transferred rate
records its donor, shift, scale and direction handling.

\subsection{Temperature, direction and isotope scaling}
Positive source rates are interpolated linearly in $\log k$ versus
$\log T$; intervals with a zero endpoint use linear interpolation in $k,T$.
Outside each usable source grid,
\begin{equation}
k(T)=k(T_e)\sqrt{T/T_e}.
\end{equation}
This supplies values on the released 10--400 K grid. It does not prescribe
a consumer's extrapolation outside the released grid.

When a copied or donor transition is used in the reverse of its source
energy ordering, source energies and weights first give
\begin{equation}
k_{lu}^{\rm source}=k_{ul}^{\rm source}
\frac{g_u^{\rm source}}{g_l^{\rm source}}
\exp[-c_2\Delta E_{\rm source}/T].
\end{equation}
The resulting coefficient is attached to the downward target transition.
Its upward target partner must then be computed with the \emph{target}
energy gap and weights. Mixing the two gaps violates the intended model.

Carbon-13 rates are transferred by identity with
\begin{equation}
k_{13,c}=\sqrt{\mu_{12,c}/\mu_{13,c}}\,k_{12,c},\qquad
\mu_{i,c}=\frac{m_i m_c}{m_i+m_c}.
\end{equation}
The isotope mass sums in the scaling module, rather than the rounded display
mass in the DAT header, define these factors. The frozen headers retain
32.04186 u for ordinary methanol (average mass convention) and
33.02957 u for carbon-13 methanol (monoisotopic convention). This
historical display inconsistency is retained to preserve exact v3 identity;
it does not enter the reduced-mass rate calculation. This assumes unchanged
state-to-state cross-section integrals at the same temperature.

\section{File use and reproducible workflow}
The literal \texttt{[EXPERIMENTAL]} tag in the frozen DAT titles is historical
and retained solely for byte-identical v3 reproduction. The maintainer's
6 November 2026 adoption supersedes that label; current delivery manifests
identify the adopted-use status. Programs displaying the raw title may still
show the old tag.

The output follows the
\href{https://home.strw.leidenuniv.nl/~moldata/molformat.html}{LAMDA format}:
level energies in cm$^{-1}$, $A$ in s$^{-1}$, line frequency in GHz and
downward collision coefficients in cm$^3$\,s$^{-1}$. Partner IDs are 2
(para-H$_2$), 3 (ortho-H$_2$) and 6 (He). The format does not fix densities,
an H$_2$ mixture, a He abundance or a fallback policy. A consumer that scales
He rows to represent missing H$_2$ rates implements an additional model;
that choice must not be confused with adding physical He collisions.

\begin{table}[ht]
\centering\small
\caption{Measured collision coverage. Each cell gives rows / levels appearing
in rows. Present zero coefficients count as rows; these are not counts of
positive couplings or of complete networks.}
\label{tab:coverage}
\begin{tabular}{lrrr}
\toprule Product & para-H$_2$ & ortho-H$_2$ & He\\\midrule
A-12 & 1,528,128 / 7,600 & 31,350 / 836 & 332,615 / 2,684\\
E-12 & 1,615,634 / 7,669 & 31,350 / 836 & 335,378 / 2,696\\
A-13 & 1,556,888 / 7,714 & 31,350 / 836 & 332,615 / 2,684\\
E-13 & 1,649,647 / 7,794 & 31,350 / 836 & 335,378 / 2,696\\\bottomrule
\end{tabular}
\end{table}
Only para-H$_2$ receives the J/K extrapolation. Ortho-H$_2$ leaves
6,764--6,958 levels per file without a row; He leaves 4,916--5,098.
For ortho fraction $f={\rm OPR}/(1+{\rm OPR})$, the file-level mixture is
\begin{equation}
C_{ul}=n({\rm H}_2)[(1-f)k_{ul,p}+f k_{ul,o}]+n({\rm He})k_{ul,\mathrm{He}}.
\end{equation}
An absent entry contributes no term, but represents unknown coupling rather
than an established physical zero. Increasing the ortho fraction can remove
much of the represented high-level H$_2$ excitation. No automatic substitution
of para for ortho or He for H$_2$ is supplied. Any such fallback is an additional
consumer model; declare it together with densities and OPR and avoid counting
physical He twice. Per-manifold coverage is in
\path{provenance/collider-coverage.json}.

The portable release contains Python source, native prediction drivers,
parity/dipole overlays, the input lock, numerical reference checksums and
tests. The four species archives can be used directly. To regenerate them, obtain
the 97 locked inputs (153,588,511 bytes) separately. These include 58 original
BELGI Fortran files. The private reproduction bundle contains them beneath
\path{inputs/}. Its source and input checksums are in \path{SOURCE-MANIFEST.json}.

\begin{lstlisting}
python scripts/l33t.py check-inputs --inputs inputs
python scripts/l33t.py rebuild --inputs inputs --output work/rebuild
python scripts/l33t.py extract --products work/rebuild/products \
  --isotope both --species both --gzip --output work/delivery
\end{lstlisting}

The rebuild performs normalization, fresh native calculations for both isotope
Hamiltonians, all radiation and collision assembly, serialization, exact
reference checks and a separate streaming structural/format audit. Output
directories must be new. The compressed provenance tables retain physical
level energies and their basis, radiative strengths/operators/frequency policy,
and collision sources/transforms/donors. Extraction copies complete selected
products with these sidecars; it does not subset levels.

The frozen Xu observation crosswalk and empirical source transcriptions are
inputs. This package does not regenerate their historical human/source
adjudication from the original full literature record. Source constants and
model choices are pinned rather than refitted. Exact product hashes test the
preserved numerical environment; different compiler/linear-algebra stacks
may require a separately documented numerical comparison if bitwise identity
fails. \texttt{verify} checks exact DAT bytes and decompressed provenance CSV
content; \texttt{--archive-bytes} also checks the original compressed bytes.
The \texttt{compare} command matches quantum labels, endpoint pairs and
collider IDs independently of transition row order. Its tolerances are
$10^{-8}$ cm$^{-1}$ absolute plus $10^{-10}$ relative for level energies;
$10^{-9}$ absolute plus $10^{-9}$ relative for frequencies in GHz and upper
energies in K; and $10^{-7}$ relative with zero absolute tolerance for $A$
and collision coefficients. Thus zero rates remain zero. Masses, weights,
grids and identity sets match exactly. Sidecar provenance is checked by
\texttt{verify}, not by the numerical DAT comparator. A numerical pass does
not silently promote a new platform's bytes to the same frozen release.

\section{Validation, limitations and release state}
\subsection{Evidence supported by the existing audits}
The v3 acceptance record covers all four full files and all provenance rows.
Relative to v2, every previous radiative row is preserved, as are the level
and complete collision sections; only 28,347 weak $\nu_{10}$ branches were
added. Independent A/E contractions through upper $J=12$ in both isotopes
matched 13,614 v3 branches to maximum relative discrepancy
$4.1\times10^{-13}$. The full consumer comparison at 250 K checked
6,372,797 assembled collision pairs, with maximum detailed-balance residual
$3.85\times10^{-15}$. This tests consumer upward-rate arithmetic with target
energies and weights, not the physical accuracy of downward coefficients.

Against the available native $c$-operator reference, the maximum omitted
$\nu_{10}$ spontaneous decay into retained lower states is
$0.0066703\%$. This accounts for the product floor only: native threshold,
finite torsional basis, incident radiation and the physical band model remain
separate. It is not an arbitrary-radiation pumping or non-LTE convergence bound.

The RF-only $J=13$--15 holdout quantifies the donor rule. For translated
pairs, the 90th-percentile factor errors are 3.57/8.82 (A/E) at 100 K and
4.23/9.35 at 200 K. The factor metric is
$\exp(|\log(k_{\rm predicted}/k_{\rm reference})|)$.
This backtest does not validate the E hybrid or extension to $J=46$.
Extrapolated para-H$_2$ rows carry 57.8--58.6\% of the LTE-weighted downward
para-H$_2$ flux at 400 K in the audited products. Such errors are model
uncertainties, not numerical solver tolerances.

\subsection{Numerical reproduction}
\path{provenance/validation.json} separates current source-package checks from
the dated rc2 full rebuild and inherited independent audits. It records the
tested platform, exact product comparison and archive checks. A subsequent
main-isotope PERMANENT calculation through $J=46$ reproduced the native
level and line tables byte for byte. Synthetic fixtures cover the absolute Einstein-A
prefactor, copied-rate reversal, signed-K translation, observed-frequency
precedence, A parity, isotope scaling, exact-tie reorientation and E-hybrid
zeros/fallbacks. Portable checks reject changed identities, grids and rate
values while accepting specified numerical tolerance and compression changes;
patch guards also operate under optimized Python. These tests are not a replacement
for full generation or independent molecular validation.

The fresh native-INTEN/cache check covers 143,176 strengths across both
isotopes in PERMANENT and VCO\_A modes through $J=12$. All endpoint sets
agree; the maximum relative discrepancy is $3.89\times10^{-12}$, within
the stated numerical tolerance. This tests equivalent contractions of the
same repaired eigenvectors, not independent Hamiltonian physics. The direct
$c$-operator check remains inherited evidence.

The remaining scientific qualifications include high-$J$ and high-$|K|$
spectroscopy; labels near mixed resonances; collision donor errors; the E
hybrid; excited-manifold copies; sparse torsion-changing collisions;
temperature continuation; pure-axis vibrational moments; missing vibrational
couplings; and isotope transfer. A recent consumer calculation shows
observable-dependent basis and collision sensitivities. Numerical agreement
of the molecular files alone cannot establish maser gain, brightness,
saturation, source geometry or a common emitting parcel.

\clearpage
\subsection{Availability and citation}
Download the four species archives and their SHA-256 checksums from the
\href{https://github.com/blankhaar/methanol-l33t/releases/tag/v1.0}{GitHub release for v1.0}.
Each ZIP contains one LAMDA file and three compressed CSV tables.
Cite Boy Lankhaar (2026), \emph{L33T methanol molecular data for excitation
calculations}, version 1.0. Give the isotope, A/E symmetry and DAT checksum.
Credit BELGI-C$_s$ and Hougen, Kleiner and Godefroid~\cite{Hougen_1994},
the relevant Hamiltonian fits and collision calculations, and the databases
identified in the provenance. No dataset DOI has been assigned.

The project code is licensed under EUPL-1.2. Upstream BELGI material and
source data retain their respective terms. Regeneration requires the 97 locked
inputs separately. Native reproduction on a second platform remains untested.
The \href{https://github.com/blankhaar/methanol-l33t/tree/v1.0/methanol-l33t/docs}{release documentation}
includes the editable companion and bibliography.

\bigskip
\begingroup
\small
\setlength{\bibsep}{2pt}
\bibliographystyle{unsrtnat}
\bibliography{references}
\endgroup
\end{document}
